% ============================================================================
%  Curatio — Server & Model Integration Walkthrough
%  Compile with:  pdflatex "server-and-model-walkthrough.tex"  (run twice for TOC)
%  Needs: amsmath, amssymb, xcolor, tcolorbox, enumitem, booktabs, listings, hyperref
% ============================================================================
\documentclass[11pt]{article}

\usepackage[utf8]{inputenc}
\usepackage[T1]{fontenc}
\usepackage[a4paper,margin=1in]{geometry}
\usepackage{amsmath,amssymb}
\usepackage{xcolor}
\usepackage[most]{tcolorbox}
\usepackage{enumitem}
\usepackage{booktabs}
\usepackage{array}
\usepackage{fancyhdr}
\usepackage{titlesec}
\usepackage{listings}
\usepackage[hidelinks]{hyperref}

% ---- Theme colours (match the Curatio app) --------------------------------
\definecolor{forest}{HTML}{166534}
\definecolor{vibrant}{HTML}{22C55E}
\definecolor{deepforest}{HTML}{064E3B}
\definecolor{slate}{HTML}{1E293B}
\definecolor{mutel}{HTML}{475569}
\definecolor{sand}{HTML}{F4F7F5}
\definecolor{codebg}{HTML}{F1F5F9}

\hypersetup{linkcolor=forest, urlcolor=forest, citecolor=forest}

\lstset{
  basicstyle=\ttfamily\small,
  backgroundcolor=\color{codebg},
  frame=single,
  framerule=0.4pt,
  rulecolor=\color{forest!40},
  breaklines=true,
  columns=fullflexible,
  keepspaces=true,
  showstringspaces=false,
  literate={_}{\_}{1},
}

\titleformat{\section}{\Large\bfseries\color{forest}}{\thesection}{0.6em}{}
\titleformat{\subsection}{\large\bfseries\color{deepforest}}{\thesubsection}{0.6em}{}
\setlist{nosep, leftmargin=1.4em, topsep=2pt}

\newtcolorbox{idea}{
  enhanced, breakable, colback=vibrant!7, colframe=vibrant!70!black,
  boxrule=0.6pt, arc=3pt, left=8pt, right=8pt, top=6pt, bottom=6pt,
  fonttitle=\bfseries, coltitle=forest,
  title={THE IDEA}, attach title to upper={\par\smallskip}}

\newtcolorbox{flowbox}{
  enhanced, breakable, colback=slate!4, colframe=slate!55,
  boxrule=0.6pt, arc=3pt, left=8pt, right=8pt, top=6pt, bottom=6pt,
  fonttitle=\bfseries, coltitle=slate,
  title={REQUEST FLOW}, attach title to upper={\par\smallskip}}

\newtcolorbox{filebox}[1]{
  enhanced, breakable, colback=sand, colframe=forest!60,
  boxrule=0.6pt, arc=3pt, left=8pt, right=8pt, top=6pt, bottom=6pt,
  fonttitle=\bfseries\ttfamily\small, coltitle=forest,
  title={#1}, attach title to upper={\par\smallskip}}

\newtcolorbox{qabox}{
  enhanced, breakable, colback=slate!5, colframe=deepforest,
  boxrule=0.6pt, arc=3pt, left=8pt, right=8pt, top=6pt, bottom=6pt,
  fonttitle=\bfseries, coltitle=deepforest,
  title={LIKELY EXAMINER QUESTIONS}, attach title to upper={\par\smallskip}}

\newcommand{\route}[1]{{\small\texttt{\textcolor{mutel}{#1}}}}

\pagestyle{fancy}
\fancyhf{}
\fancyhead[L]{\small\textcolor{forest}{Curatio}}
\fancyhead[R]{\small\textcolor{mutel}{Server \& Model Walkthrough}}
\fancyfoot[C]{\small\thepage}
\renewcommand{\headrulewidth}{0.4pt}

% ============================================================================
\begin{document}

\begin{titlepage}
\centering
\vspace*{3cm}
{\Huge\bfseries\color{forest} Curatio}\\[0.4cm]
{\LARGE\bfseries How the Server and the BioBERT Model Work Together}\\[1cm]
{\large A plain-English guide to the API layer, the ML service, and the fine-tuned model.}\\[2cm]
{\large Final Year Project}\\[0.3cm]
{\large \today}\\[2cm]
\begin{tcolorbox}[width=0.88\textwidth, colback=sand, colframe=forest, arc=4pt]
\small\centering
This document explains the \textbf{runtime architecture}: what happens when a patient types a chief
complaint in the browser, how the Node.js API forwards it to Python, how BioBERT runs inference,
and what JSON comes back. Read this before your demo or viva if someone asks ``how does the backend
actually work?''
\end{tcolorbox}
\end{titlepage}

\tableofcontents
\newpage

% ===========================================================================
\section{The big picture in one minute}

Curatio is \textbf{not} one monolithic program. At runtime there are three cooperating pieces:

\begin{enumerate}
  \item \textbf{Client} (React/Vite) --- the web UI in the browser. It collects text and displays
        the prediction. It never loads the 414\,MB model.
  \item \textbf{API server} (Node.js/Express on port 5000) --- a thin \emph{gateway}: CORS,
        JSON validation, health checks, and forwarding requests to the ML service.
  \item \textbf{ML service} (Python/FastAPI on port 8001) --- loads BioBERT once at startup, runs
        tokenisation + inference, maps acuity to SATS colour, and flags low-confidence cases.
\end{enumerate}

\begin{idea}
\textbf{Why split Node and Python?}
\begin{itemize}
  \item BioBERT lives in the \textbf{Python/Hugging Face} ecosystem (PyTorch, \texttt{transformers}).
  \item The web stack is \textbf{JavaScript/React} --- natural for HTTP, CORS, and deployment to
        Render/Vercel.
  \item The API is a \textbf{proxy}: it keeps the browser simple (one URL) while the heavy ML work
        stays in a dedicated Python process that can be scaled or redeployed independently.
\end{itemize}
\end{idea}

\subsection{Architecture diagram}

\begin{center}
\begin{tcolorbox}[width=0.95\textwidth, colback=white, colframe=forest, arc=4pt, boxrule=0.8pt]
\small
\begin{tabular}{@{}c@{\quad$\longrightarrow$\quad}c@{\quad$\longrightarrow$\quad}c@{\quad$\longrightarrow$\quad}c@{}}
\textbf{Browser} &
\textbf{Express API} &
\textbf{FastAPI ML} &
\textbf{BioBERT weights} \\[-2pt]
{\footnotesize React SPA} &
{\footnotesize \route{curatio/server}} &
{\footnotesize \route{curatio/server/ml}} &
{\footnotesize local dir or HF Hub} \\
\end{tabular}

\vspace{0.6em}
\hrule height 0.4pt
\vspace{0.4em}

\textbf{Local dev:}\quad
\texttt{localhost:5173} $\to$ proxy \texttt{/api} $\to$ \texttt{:5000} $\to$ \texttt{:8001} $\to$ \texttt{MODEL\_PATH}

\vspace{0.3em}

\textbf{Production:}\quad
Vercel $\to$ Render $\to$ Hugging Face Space $\to$ \texttt{MODEL\_ID} on the Hub
\end{tcolorbox}
\end{center}

% ===========================================================================
\section{Starting the stack locally}

From \route{curatio/server}, one command starts \textbf{both} processes:

\begin{lstlisting}
npm run dev
\end{lstlisting}

This uses \texttt{concurrently} to run:
\begin{itemize}
  \item \textbf{api} --- \texttt{nodemon index.js} $\to$ Express on \texttt{http://localhost:5000}
  \item \textbf{ml} --- \texttt{uvicorn app:app --reload --port 8001} $\to$ FastAPI on \texttt{http://127.0.0.1:8001}
\end{itemize}

The client (\route{curatio/client}) runs separately:

\begin{lstlisting}
npm run dev    # Vite on http://localhost:5173
\end{lstlisting}

Vite proxies any path starting with \texttt{/api} to port 5000, so the browser can call
\texttt{/api/triage/predict} without CORS issues during development.

% ===========================================================================
\section{End-to-end: one prediction request}

This is the path when someone submits ``central crushing chest pain'' on the \textbf{Test BioBERT
Triage} page.

\subsection{Step 1 --- Browser sends JSON}

\begin{filebox}{Client: TriageTestPage.jsx}
\begin{lstlisting}
POST /api/triage/predict
Content-Type: application/json

{ "text": "central crushing chest pain, radiating to left arm" }
\end{lstlisting}
\end{filebox}

\begin{itemize}
  \item Locally, \texttt{VITE\_API\_URL} is empty, so the URL is relative and hits the Vite proxy.
  \item In production, \texttt{VITE\_API\_URL} is set to the Render API origin (e.g.\
        \texttt{https://curatio-api.onrender.com}).
\end{itemize}

\subsection{Step 2 --- Express validates and forwards}

\begin{filebox}{Server: routes/triage.js}
Express receives the body, checks that \texttt{text} is a non-empty string, then \textbf{does not}
run any ML itself. It forwards the same JSON to the Python service:

\begin{lstlisting}
POST ${ML_SERVICE_URL}/predict
{ "text": "central crushing chest pain, radiating to left arm" }
\end{lstlisting}

\texttt{ML\_SERVICE\_URL} defaults to \texttt{http://127.0.0.1:8001} locally; in production it
points to your Hugging Face Space URL.
\end{filebox}

The Express layer also:
\begin{itemize}
  \item Applies \textbf{CORS} --- only allows origins listed in \texttt{CLIENT\_URL} (plus localhost
        in dev).
  \item Exposes \texttt{GET /api/health} (API alive) and \texttt{GET /api/triage/health} (proxies
        to ML \texttt{/health}).
  \item Returns \texttt{503} with a helpful message if Python is not running.
\end{itemize}

\subsection{Step 3 --- FastAPI receives and calls \texttt{predict()}}

\begin{filebox}{ML service: ml/app.py}
On startup, the \texttt{lifespan} hook calls \texttt{load\_model()} once --- BioBERT is loaded into
memory before the first request.

\begin{lstlisting}
@app.post("/predict")
def predict_endpoint(body: PredictRequest):
    return predict(body.text)
\end{lstlisting}

\texttt{PredictRequest} enforces \texttt{text} length 1--2000 characters (Pydantic validation).
\end{filebox}

\subsection{Step 4 --- Inside \texttt{predict()}: what the model actually does}

\begin{filebox}{ML service: ml/predict.py}
\begin{enumerate}
  \item \textbf{Load pipeline} (singleton) --- first call loads tokenizer + model; later calls reuse
        the in-memory pipeline.
  \item \textbf{Tokenise} --- \texttt{pipe(text, truncation=True, max\_length=128)}. Only the first
        128 tokens are kept (matches training and presentation slides).
  \item \textbf{Forward pass} --- 12 transformer layers produce logits for 5 classes
        (\texttt{LABEL\_0} \ldots \texttt{LABEL\_4}).
  \item \textbf{Softmax} --- Hugging Face \texttt{text-classification} pipeline returns a score
        (probability) for every class.
  \item \textbf{Map labels to acuity} --- \texttt{LABEL\_k} $\to$ acuity level \texttt{k+1} (levels
        1--5).
  \item \textbf{Pick winner} --- highest probability = predicted acuity; that value = \textbf{confidence}.
  \item \textbf{SATS mapping} --- level 1 $\to$ Red, 2 $\to$ Orange, 3 $\to$ Yellow, 4--5 $\to$ Green.
  \item \textbf{Confidence gate} --- if \texttt{confidence < CONFIDENCE\_THRESHOLD} (default 0.85),
        set \texttt{bayesian\_candidate: true}.
\end{enumerate}
\end{filebox}

\begin{idea}
This is \textbf{only the NLP layer} from your presentation (Step 4 / BioBERT). The live test page
returns BioBERT's language-based triage. Full Curatio fusion (TEWS + Bayesian + safety rules) is
the \emph{research design} shown in slides; the deployed test endpoint implements the
\textbf{fine-tuned classifier} part so you can demo real inference.
\end{idea}

\subsection{Step 5 --- JSON response travels back}

Python returns a dict $\to$ Express passes it through unchanged $\to$ React displays it.

Example response:

\begin{lstlisting}
{
  "text": "central crushing chest pain, radiating to left arm",
  "predicted_acuity_level": 2,
  "predicted_class_index": 1,
  "confidence": 0.94,
  "probabilities": [
    { "level": 2, "class_index": 1, "probability": 0.94 },
    { "level": 3, "class_index": 2, "probability": 0.04 },
    ...
  ],
  "sats_colour": "Orange",
  "bayesian_candidate": false
}
\end{lstlisting}

\begin{itemize}
  \item \texttt{predicted\_acuity\_level} --- the winning class (1 = most urgent).
  \item \texttt{confidence} --- $\max_c \hat{y}_c$ from the presentation (softmax peak).
  \item \texttt{bayesian\_candidate} --- \texttt{true} when confidence $< 0.85$; in the full system
        this would trigger the Bayesian fallback at fusion.
  \item \texttt{probabilities} --- full distribution for transparency / debugging.
\end{itemize}

% ===========================================================================
\section{Where the model weights live}

The ML service resolves the model source at startup:

\begin{center}
\begin{tabular}{@{}lll@{}}
\toprule
\textbf{Environment} & \textbf{Variable} & \textbf{Behaviour} \\
\midrule
Local dev & \texttt{MODEL\_PATH} & Load from disk under \route{fine-tuned-biobert/...} \\
Production & \texttt{MODEL\_ID} & Download from Hugging Face Hub at startup \\
Fallback & (default path) & Same local folder if \texttt{MODEL\_PATH} unset \\
\bottomrule
\end{tabular}
\end{center}

\begin{filebox}{predict.py --- resolve\_model\_source()}
\begin{lstlisting}
if MODEL_ID is set:
    source = "your-username/curatio-biobert-triage"   # Hub repo
else:
    source = local MODEL_PATH directory               # model.safetensors on disk
\end{lstlisting}
\end{filebox}

Weights are loaded via Hugging Face \texttt{AutoTokenizer} and
\texttt{AutoModelForSequenceClassification.from\_pretrained(source)}. The service builds a
\texttt{text-classification} pipeline with \texttt{top\_k=None} so \textbf{all five} class
probabilities are returned (needed for confidence and for demo tables).

Device: CPU by default (\texttt{device=-1}); uses CUDA if PyTorch sees a GPU.

% ===========================================================================
\section{File map: who owns what}

\begin{center}
\begin{tabular}{@{}p{0.32\textwidth}p{0.58\textwidth}@{}}
\toprule
\textbf{File} & \textbf{Responsibility} \\
\midrule
\route{server/index.js} &
Express app: CORS, JSON middleware, mounts \texttt{/api/triage}, \texttt{/api/health}. \\[4pt]
\route{server/routes/triage.js} &
Proxy routes: \texttt{POST /predict}, \texttt{GET /health} $\to$ \texttt{ML\_SERVICE\_URL}. \\[4pt]
\route{server/ml/app.py} &
FastAPI app: lifespan loads model; HTTP endpoints \texttt{/health}, \texttt{/predict}. \\[4pt]
\route{server/ml/predict.py} &
Core logic: load model, run pipeline, map acuity/SATS, confidence gate. \\[4pt]
\route{server/ml/Dockerfile} &
Container for Hugging Face Spaces (uvicorn on port 7860). \\[4pt]
\route{client/vite.config.js} &
Dev proxy: \texttt{/api} $\to$ \texttt{localhost:5000}. \\[4pt]
\route{client/.../TriageTestPage.jsx} &
UI that calls \texttt{/api/triage/predict} and renders probabilities. \\
\bottomrule
\end{tabular}
\end{center}

% ===========================================================================
\section{Health checks and failure modes}

\subsection{Health endpoints}

\begin{itemize}
  \item \texttt{GET /api/health} --- Express only: \texttt{\{ status: "ok", service: "curatio-server" \}}.
  \item \texttt{GET /api/triage/health} --- Express forwards to ML \texttt{/health}.
  \item \texttt{GET /health} (ML) --- reports \texttt{model\_path}, \texttt{model\_source},
        \texttt{weights\_found}, \texttt{model\_loaded}.
\end{itemize}

Use these before a live demo to warm cold-started services (Render free tier, HF Spaces sleep).

\subsection{Common errors}

\begin{center}
\begin{tabular}{@{}llp{0.42\textwidth}@{}}
\toprule
\textbf{HTTP} & \textbf{Cause} & \textbf{Fix} \\
\midrule
400 & Empty \texttt{text} & Client validation / user input \\
503 & ML not running & \texttt{npm run dev} from \route{curatio/server} \\
503 & Weights missing & Set \texttt{MODEL\_PATH} or upload Hub + \texttt{MODEL\_ID} \\
CORS error & Wrong origin & Set \texttt{CLIENT\_URL} on Render to Vercel URL \\
Slow first request & Cold start & Hit \texttt{/health} 1--2 min before demo \\
\bottomrule
\end{tabular}
\end{center}

% ===========================================================================
\section{Production deployment (how the pieces connect)}

Deploy order: \textbf{Hub model $\to$ HF Space $\to$ Render API $\to$ Vercel client}.

\begin{enumerate}
  \item Upload \texttt{model.safetensors} to Hugging Face Hub $\to$ \texttt{MODEL\_ID}.
  \item Deploy \route{server/ml/} as a Docker Space $\to$ \texttt{ML\_SERVICE\_URL}.
  \item Deploy \route{server/} to Render with \texttt{ML\_SERVICE\_URL} and \texttt{CLIENT\_URL}.
  \item Deploy \route{client/} to Vercel with \texttt{VITE\_API\_URL} = Render URL.
\end{enumerate}

Full step-by-step commands live in \route{curatio/DEPLOYMENT.md}.

\begin{idea}
In production the browser \textbf{never} talks to Python directly. It only sees the Render API.
That keeps the ML service private, lets you rotate models without redeploying the client, and
matches standard microservice practice.
\end{idea}

% ===========================================================================
\section{How this connects to your presentation maths}

\begin{center}
\begin{tabular}{@{}p{0.28\textwidth}p{0.62\textwidth}@{}}
\toprule
\textbf{Presentation concept} & \textbf{Where it happens in code} \\
\midrule
Token limit 128 &
\texttt{predict.py}: \texttt{max\_length=128}, \texttt{truncation=True} \\[4pt]
12 encoder layers &
Inside loaded BioBERT weights (not reimplemented in app code) \\[4pt]
Softmax over 5 levels &
Hugging Face pipeline; scores in \texttt{probabilities[]} \\[4pt]
Confidence $\max \hat{y}_c$ &
\texttt{confidence = best["probability"]} \\[4pt]
Gate if $< \tau$ &
\texttt{bayesian\_candidate = confidence < CONFIDENCE\_THRESHOLD} \\[4pt]
Acuity $\to$ SATS colour &
\texttt{SATS\_BY\_ACUITY} dict in \texttt{predict.py} \\[4pt]
TEWS / Bayesian / Fusion &
Designed in slides; not in this HTTP predict path (future full pipeline) \\
\bottomrule
\end{tabular}
\end{center}

When an examiner asks ``does the server \emph{implement} fusion?'', the honest answer is:
\textbf{the deployed predict endpoint runs BioBERT inference plus a confidence flag}. Fusion,
TEWS, and Bayesian are documented and demonstrated in the presentation; extending the API to run
the full \texttt{f\_fusion} controller would be the next engineering step.

% ===========================================================================
\section{Environment variables cheat sheet}

\begin{center}
\small
\begin{tabular}{@{}llll@{}}
\toprule
\textbf{Variable} & \textbf{Set on} & \textbf{Example} & \textbf{Purpose} \\
\midrule
\texttt{MODEL\_PATH} & ML (local) & path to fine-tuned folder & Load weights from disk \\
\texttt{MODEL\_ID} & ML (prod) & \texttt{user/curatio-biobert-triage} & Pull from Hub \\
\texttt{ML\_SERVICE\_URL} & Express & \texttt{http://127.0.0.1:8001} & Where to forward predict \\
\texttt{PORT} & Express & 5000 (Render injects) & API listen port \\
\texttt{CLIENT\_URL} & Express & Vercel URL & CORS allow-list \\
\texttt{CONFIDENCE\_THRESHOLD} & ML (+ Express copy) & 0.85 & Bayesian candidate flag \\
\texttt{VITE\_API\_URL} & Client build & Render URL & Production API origin \\
\bottomrule
\end{tabular}
\end{center}

% ===========================================================================
\section{What to say in the viva (30-second version)}

\begin{quote}
\emph{``When a nurse or patient submits a chief complaint, the React client POSTs JSON to our
Express API on Render. Express doesn't run the model --- it validates the text and forwards the
request to a separate Python FastAPI service that loads our fine-tuned BioBERT once at startup.
The model tokenises the text to 128 tokens, runs a forward pass through twelve transformer layers,
applies softmax over five acuity classes, and returns the predicted level, full probabilities, a
SATS colour mapping, and a flag if confidence is below 0.85. That flag corresponds to the Bayesian
fallback in our fusion design. Locally, both processes start with one \texttt{npm run dev}; in
production the model weights live on the Hugging Face Hub so we don't bundle 414 megabytes in the
API container.''}
\end{quote}

% ===========================================================================
\section{Likely examiner questions}

\begin{qabox}
\textbf{Q: Why not run BioBERT inside Node.js?}\\
A: The model is PyTorch + Hugging Face \texttt{transformers}. Python is the standard runtime for
that stack. Node handles HTTP/CORS/deployment ergonomics; Python handles inference. They communicate
over a simple REST JSON contract.

\medskip
\textbf{Q: What if the ML service crashes?}\\
A: Express catches the connection error and returns HTTP 503 with ``ML service unavailable'' so the
UI can show a clear message instead of hanging.

\medskip
\textbf{Q: Is the model loaded on every request?}\\
A: No. \texttt{load\_model()} uses a module-level singleton (\texttt{\_pipeline}). FastAPI's
lifespan loads it at startup; subsequent requests reuse the same in-memory pipeline.

\medskip
\textbf{Q: How do you know the server uses \emph{your} fine-tuned weights?}\\
A: \texttt{MODEL\_PATH} or \texttt{MODEL\_ID} points explicitly to the fine-tuned checkpoint
(trained on $\sim$80k merged chief complaints). The \texttt{/health} endpoint reports which source
was loaded.

\medskip
\textbf{Q: Does the API implement TEWS or Bayesian?}\\
A: Not in the current \texttt{/predict} endpoint --- that runs the BioBERT NLP layer and sets
\texttt{bayesian\_candidate} when confidence is low. TEWS and Bayesian fusion are specified in the
system design and presentation; the test page demonstrates the language model path end-to-end.
\end{qabox}

% ===========================================================================
\section{Quick reference: sequence diagram}

\begin{flowbox}
\begin{enumerate}
  \item User clicks \textbf{Predict} in browser.
  \item \texttt{fetch("/api/triage/predict", \{ text \})}
  \item Vite proxy (dev) or Vercel $\to$ Render (prod)
  \item Express \texttt{routes/triage.js} validates \texttt{text}
  \item Express \texttt{fetch(ML\_SERVICE\_URL + "/predict")}
  \item FastAPI \texttt{predict\_endpoint} $\to$ \texttt{predict(text)}
  \item BioBERT pipeline: tokenise $\to$ encode $\to$ softmax $\to$ map acuity/SATS
  \item JSON response $\uparrow$ Express $\uparrow$ browser $\to$ UI shows level, colour, bars
\end{enumerate}
\end{flowbox}

\vfill
\noindent\rule{\textwidth}{0.4pt}\\[4pt]
{\small\textcolor{mutel}{See also: \route{curatio/DEPLOYMENT.md} (deploy steps),
\route{presentation-walkthrough.tex} (slide-by-slide defence script). Compile this file with
\texttt{pdflatex} twice for the table of contents.}}

\end{document}
